\chapter{Couche JSF : beans, portées et navigation}\label{chap:jsf}

\section{Cycle de vie et conséquences}
Une requête JSF traverse six phases : restauration de la vue, application des valeurs de requête, validation,
mise à jour du modèle, invocation de l'application, rendu. Deux conséquences structurent tout le code.

\begin{itemize}
  \item Une valeur saisie n'arrive dans le bean qu'à la phase de \textit{mise à jour du modèle}, et seulement si la
        validation a réussi. Le bean ne lit donc jamais \code{request.getParameter} : il reçoit des propriétés.
  \item Une méthode d'action (\code{action="\#\{bean.methode\}"}) s'exécute à la phase d'\textit{invocation de
        l'application}, après la mise à jour du modèle. Elle retourne un résultat de navigation.
\end{itemize}

\section{Choix des portées}
La portée d'un bean décide de la durée de vie de son état. Le tableau~\ref{tab:portees} justifie chaque choix.

\begin{table}[htbp]
\centering
\caption{Portées des beans}\label{tab:portees}
\small
\begin{tabularx}{\linewidth}{@{}llX@{}}
\toprule
\textbf{Bean} & \textbf{Portée} & \textbf{Justification}\\
\midrule
\code{ConnexionBean}, \code{InscriptionBean} & \code{@ViewScoped} & Garde la saisie tant que la même page est réaffichée (erreur de validation) ; disparaît à la navigation.\\
\code{VitrineBean}, \code{DetailBean} & \code{@ViewScoped} & Les filtres arrivent par l'adresse (paramètres de vue) ; l'état n'a pas à survivre à la page.\\
\code{PanierBean} & \code{@ViewScoped} & Le panier lui-même est lu en base à chaque affichage ; le bean ne garde que les quantités saisies.\\
\code{ProduitAdminBean}, \code{CategorieAdminBean} & \code{@ViewScoped} & Un formulaire d'édition par page.\\
\code{UserSession} & \code{@RequestScoped} & Façade sans état propre : l'état de l'utilisateur est dans la \code{HttpSession}, sous forme d'un objet \code{SessionUser} sérialisable.\\
\code{LoginGuard}, \code{Affichage} & \code{@ApplicationScoped} & Partagés : le compteur de tentatives doit voir tous les visiteurs, et les aides d'affichage n'ont pas d'état.\\
\bottomrule
\end{tabularx}
\end{table}

Deux décisions méritent d'être détaillées.

\textbf{Aucun bean en \code{@SessionScoped}.} Une portée session est tentante pour « se souvenir de l'utilisateur »,
mais elle fait vivre l'état d'un bean aussi longtemps que la session. Ici, la session ne contient qu'un petit
objet immuable (\code{SessionUser} : identifiant, nom, rôle) ; \code{UserSession} est une façade de requête qui le lit et
le modifie. La fixation de session se traite ainsi à un seul endroit (section~\ref{sec:usersession}).

\textbf{\code{@ViewScoped} de JSF, pas de CDI.} L'annotation utilisée est \code{javax.faces.view.ViewScoped}, couplée à
\code{@Named}. Les beans de cette portée sont placés dans l'état de la vue : ils doivent être \code{Serializable}, ce qui a conduit à
rendre sérialisables les objets formulaires (\code{ProduitForm}, \code{CategorieForm}, \code{InscriptionForm}) et à
remplacer les \code{Map} génériques par des \code{LinkedHashMap}.

\section{Navigation et motif Post/Redirect/Get}
Une action qui modifie des données doit se terminer par une redirection, sinon un rechargement de la page
(\code{F5}) rejoue le formulaire \cite{prg}. JSF y répond par la navigation implicite : une méthode retourne
le nom d'une page suivi de \code{?faces-redirect=true}. Les paramètres utiles (\code{includeViewParams=true}) sont
conservés dans l'adresse.

Le message de confirmation pose un problème : il est créé pendant la requête qui poste, mais affiché dans la page
suivante. La solution est \code{FacesMessage} conservé par le \textit{flash} de JSF. La classe utilitaire
\code{Messages} regroupe les trois niveaux (succès, information, erreur) et la méthode \code{apresRedirection()}
qui active la conservation.

\lstinputlisting[language=Java,caption={\texttt{PanierBean} : toute action passe par une méthode unique},label={lst:executer}]{extraits/PanierBean-executer.txt}

La méthode \code{executer} évite de répéter quatre fois le même bloc. Elle capture l'exception métier, la
transforme en message rouge, vide les valeurs mises en cache (\code{panier = null}) pour que l'écran suivant relise la
base, active la conservation du message, et retourne la cible. Le motif est le même dans tous les beans de gestion.

\section{Paramètres de vue}
La vitrine se filtre par catégorie, mot-clé et page. Ces trois valeurs sont portées par l'adresse, pas par un
formulaire JSF : la recherche est un simple formulaire \code{GET}. Les paramètres sont déclarés dans la page avec
\code{f:viewParam}, et \code{f:viewAction} appelle \code{init()} une fois qu'ils sont injectés.

\lstinputlisting[language=XML,caption={Métadonnées de \texttt{vitrine.xhtml}},label={lst:meta}]{extraits/vitrine-metadata.txt}

L'avantage est concret : une recherche se partage et se met en favori (\code{vitrine.xhtml?categorie=2\&q=cle}), le bouton
« page précédente » du navigateur fonctionne, et la pagination se construit avec de simples liens. La construction de
ces adresses est un calcul pur (\code{VitrineBean.construireUrl}), vérifié par des tests sans serveur : les
paramètres sont encodés en UTF-8, la page~1 est omise et un \code{\&} saisi dans la recherche ne peut pas injecter un
paramètre supplémentaire.

\section{Session et authentification}\label{sec:usersession}
\code{UserSession.ouvrir} est le seul endroit où une connexion réussie est enregistrée. Il enchaîne quatre opérations dans cet
ordre précis.

\lstinputlisting[language=Java,caption={Ouverture de session},label={lst:ouvrir}]{extraits/UserSession-ouvrir.txt}

\begin{enumerate}
  \item Lire l'adresse de retour mémorisée par le filtre d'authentification, avant de changer d'identifiant.
  \item \code{changeSessionId()} : un nouvel identifiant de session est émis (défense contre la fixation de session).
  \item Enregistrer le \code{SessionUser} dans la nouvelle session.
  \item Renouveler le jeton anti-CSRF, lié à l'identité qui vient de s'authentifier.
\end{enumerate}

\code{ConnexionBean.connecter} appelle cette méthode après avoir interrogé le service et le \code{LoginGuard}, qui limite
les essais à cinq par quinze minutes pour un couple adresse IP / adresse électronique. Le mot de passe saisi est vidé
dans tous les cas d'échec, pour qu'il ne soit jamais renvoyé à la page.

\lstinputlisting[language=Java,caption={\texttt{ConnexionBean.connecter}},label={lst:connecter}]{extraits/ConnexionBean-connecter.txt}

\section{Aides d'affichage dans Java}
Le bean \code{Affichage} (portée application) fournit le format des prix (deux décimales, virgule, espace insécable,
MAD), les initiales d'un nom, les classes CSS de stock, le libellé du bouton d'ajout et les textes de la page d'erreur.
Ces règles auraient pu s'écrire en EL dans les pages. Elles sont en Java pour une raison pratique expliquée au
chapitre~\ref{chap:incidents} : les analyseurs statiques d'Eclipse signalent comme erreurs des expressions
conditionnelles parfaitement valides à l'exécution. Déplacer la logique dans Java supprime ces faux positifs et la rend
testable (section~\ref{sec:webchecks}).
